\documentclass[10pt]{article}

\usepackage[margin=0.75in]{geometry}
\usepackage{amsmath,amsthm,amssymb,bm}
\usepackage{xcolor}
\usepackage{cancel}
\usepackage{graphicx}
\usepackage{changepage}
\usepackage{circuitikz}
\usepackage{pgfplots}
\usepackage{physics}
\usepackage{hyperref}
\usepackage{siunitx}
\usepackage{minted}
\usepackage{fontspec}
\usepackage{relsize}
\usepackage{subfig}
\usepackage{mhchem}
\usepackage{todonotes}
\usepackage{biblatex}
\usepackage{multicol, multirow, booktabs}
\usepackage[breakable]{tcolorbox}
\usepackage[inline]{enumitem}
\patchcmd{\thebibliography}{\section*{\refname}}{}{}{}
\addbibresource{sources.bib}

\theoremstyle{definition}
\newtheorem{problem}{Problem}
\newtheorem{soln}{Solution}

\pgfplotsset{compat=newest}
\usetikzlibrary{lindenmayersystems}
\usetikzlibrary{arrows}
\usetikzlibrary{calc}
\usetikzlibrary{positioning, fit}
\usetikzlibrary{3d, perspective}

\definecolor{incolor}{HTML}{303F9F}
\definecolor{outcolor}{HTML}{D84315}
\definecolor{cellborder}{HTML}{CFCFCF}
\definecolor{cellbackground}{HTML}{F7F7F7}
\newcommand{\ui}{\hat{i}}
\newcommand{\uj}{\hat{j}}
\newcommand{\uk}{\hat{k}}
\newcommand{\ux}{\hat{x}}
\newcommand{\uy}{\hat{y}}
\newcommand{\uz}{\hat{z}}
\newcommand{\uv}[1]{\hat{\bm{{#1}}}}
\newcommand{\pr}[1]{{#1^\prime}}
\newcommand{\justif}[2]{&{#1}&\text{#2}}
\newcommand{\dul}[1]{\underline{\underline{#1}}}
\newcommand{\pdiff}[2][]{{\frac{\partial^{#1}}{\partial {{#2}^{#1}}}}}
% \newcommand{\dif}[2][]{{\frac{d^{#1}}{d{{#2}^{#1}}}}}
\pgfdeclarelayer{background}  
\pgfsetlayers{background,main}
\AtBeginDocument{\RenewCommandCopy\qty\SI}

\makeatletter
\newcommand{\boxspacing}{\kern\kvtcb@left@rule\kern\kvtcb@boxsep}
\makeatother
\newcommand{\prompt}[4]{
    \ttfamily\llap{{\color{#2}[#3]:\hspace{3pt}#4}}\vspace{-\baselineskip}
}

\newcommand{\thevenin}[2]{
  \begin{center}
    \begin{circuitikz} \draw
      (0,0) -- (2,0) to[battery1, l_=$V_{Th}\eq#1$] (2,2) 
      to[resistor, l_=$R_{Th}\eq#2$] (0,2)
      ;
      \draw [o-] (-.07,2.079);
      \draw [o-] (-.07,0.079);
    \end{circuitikz}
  \end{center}
}

\newcommand{\norton}[2]{
  \begin{center}
    \begin{circuitikz} \draw
      (0,0) -- (3,0) to[american current source, l_=$I_{N}\eq#1$] (3,2) -- (0,2) (2,0)
      to[resistor, l=$R_{N}\eq#2$] (2,2)
      ;
      \draw [o-] (-.07,2.079);
      \draw [o-] (-.07,0.079);
    \end{circuitikz}
  \end{center}
}

\newcommand{\highlight}[1]{\colorbox{yellow}{$\displaystyle #1$}}

\newcommand{\ti}[1]{\widetilde{#1}}

\newfontface{\Kaufmann}{Kaufmann}
\DeclareTextFontCommand{\kf}{\Kaufmann}
\newcommand{\scriptr}{\fontsize{12pt}{12pt}\kf{r}}

\newfontface{\KaufmannB}{Kaufmann Bold}
\DeclareTextFontCommand{\kfb}{\KaufmannB}
\newcommand{\bscriptr}{\fontsize{12pt}{12pt}\kfb{r}}

\newcommand{\bv}[1]{\vec{#1}}

\title{Physics 4605H: Assignment V}
\author{Jeremy Favro (0805980) \\ Trent University, Peterborough, ON, Canada}
\date{\today}

\begin{document}
\maketitle

% PROBLEM 2
\setcounter{problem}{1}
\setcounter{soln}{1}
\begin{problem}
~
\begin{enumerate}[label=(\alph*)]
  \item Write out all the terms of $$\left(1+x/3\right)^3.$$
  \item Write out $$(1+x/N)^N$$ for general $N$ showing up to $x^3$ terms.
  \item Write $e^x$ as an expansion in powers, showing up to the $x^3$ term. As $N\to\infty$ in (b), do your answer to (b) and (c) match?
  \item Write out all the terms of $$\prod_{i=1}^3(1+x_i).$$
  \item What is $$\prod_{i=1}^N(1+x_i)$$ for general $N$ showing up to 2nd order terms?
  \item Write $$e^{\sum_{i}x_i}$$
        as an expansion, showing up to second order terms. As $N\to\infty$ in (e), do your answers to (e) and (f) match?
\end{enumerate}
\end{problem}
\begin{soln}
  \begin{enumerate}[label=(\alph*)]
    \item \begin{align*}
            \left(1+x/3\right)^3 & =\left(1+x/3\right)\left(1+x/3\right)\left(1+x/3\right)                  \\
                                 & =\left(1+x/3+x/3+x^2/9\right)\left(1+x/3\right)                          \\
                                 & =1+x/3+x/3+x^2/9+x/3+x^2/9+x^2/9+x^3/27=1+x+\frac{x^2}{3}+\frac{x^3}{27}
          \end{align*}
    \item The general expansion for such a binomial is
          $$(1+x/N)^N=\sum_{k=0}^N\begin{pmatrix}
              N \\
              k
            \end{pmatrix}\left(\frac{x}{N}\right)^k$$
          which, up to $x^3$ expands as
          $$
            \begin{pmatrix}
              N \\
              0
            \end{pmatrix}\left(\frac{x}{N}\right)^0
            +
            \begin{pmatrix}
              N \\
              1
            \end{pmatrix}\left(\frac{x}{N}\right)^1+
            \begin{pmatrix}
              N \\
              2
            \end{pmatrix}\left(\frac{x}{N}\right)^2+
            \begin{pmatrix}
              N \\
              3
            \end{pmatrix}\left(\frac{x}{N}\right)^3
          $$
          which simpifies to
          $$
            1
            +
            \begin{pmatrix}
              N \\
              1
            \end{pmatrix}\frac{x}{N}+
            \begin{pmatrix}
              N \\
              2
            \end{pmatrix}\frac{x^2}{N^2}+
            \begin{pmatrix}
              N \\
              3
            \end{pmatrix}\frac{x^3}{N^3}
          $$
    \item Using Taylor series about $x=0$ we obtain
          $$e^x=\sum_{n=0}^3\frac{e^0}{n!}x^n=1+x+\frac{x^2}{2}+\frac{x^3}{6}$$
          because $\frac{d}{dx}e^x=e^x$. Noting that
          $$\begin{pmatrix}
              N \\
              k
            \end{pmatrix}=\frac{N!}{(N-k)!k!}$$
          which for large $N$, which we have, specifcally $N\gg k$ becomes approximately
          $$\frac{N!}{N!k!}=\frac{1}{k!}$$ which recovers our result in (c).
    \item We get
          $$
            (1+x_1)(1+x_2)(1+x_3)=1+x_1+x_2+x_3+x_1x_2+x_1x_3+x_2x_3+x_1x_2x_3
          $$
          provided whatever $x_i$ comes from is commutative.
    \item I spent a while trying to figure out a general form for this but couldn't figure it out, so I'll just write the general terms,
          $$1+x_1+x_2+\dots+x_N+x_1x_2+x_1x_3+\dots+x_1x_N+\dots x_2x_3+\dots x_2x_N+\dots x_{N-1}x_N$$
    \item \begin{align*}
            e^{\sum_{i}x_i} & =e^{x_1+x_2+x_3+\dots}      \\
                            & =e^{x_1}e^{x_2}e^{x_3}\dots \\
                            & =\left[\sum_{n=0}^{N}\frac{x_1^n}{n!}\right]\left[\sum_{n=0}^{N}\frac{x_2^n}{n!}\right]\left[\sum_{n=0}^{N}\frac{x_3^n}{n!}\right]\dots \\
                            & =\left(1+x_1+\frac{x_1^2}{2!}+\frac{x_1^3}{3!}\right)\left(1+x_2+\frac{x_2^2}{2!}+\frac{x_2^3}{3!}\right)\left(1+x_3+\frac{x_3^2}{2!}+\frac{x_3^3}{3!}\right)\dots\left(1+x_N+\frac{x_N^2}{2!}+\frac{x_N^3}{3!}\right) \\
                            & =\left(1+x_1+x_2+x_3+x_1x_2+x_1x_3+\dots x_1x_N+\dots\right)+\text{higher order terms} \\
          \end{align*}
          which agrees with my answer in (e).
  \end{enumerate}
\end{soln}

% PROBLEM 3
\begin{problem}
The density matrix at time $t$ can be expressed as
$$\hat{\rho}(t)=\sum_ip_i\ket{\psi_i(t)}\bra{\psi_i(t)}$$
where $p_i$ are probabilities such that $\sum_ip_i=1$.
$\left\{\ket{\psi_i(t)}\right\}$ are a set of pure quantum states which satisfy the time-dependent Schr\"odinger equation,
$$i\hbar\frac{\partial}{\partial t}\ket{\psi_i(t)}=\hat{H}\ket{\psi_i(t)}.$$
Show that
$$i\hbar \frac{\partial}{\partial t}\hat{\rho}=\left[\hat{H},\hat{\rho}(t)\right]$$
\end{problem}
\begin{soln}
  Okay,
  \begin{align*}
    i\hbar \frac{\partial}{\partial t}\hat{\rho} & =i\hbar \pdiff{t}\sum_ip_i\ket{\psi_i(t)}\bra{\psi_i(t)}                                                                                           \\
                                                 & =i\hbar \sum_ip_i\pdiff{t}\ket{\psi_i(t)}\bra{\psi_i(t)}\justif{\quad}{The derivative is a linear operator}                                        \\
                                                 & =i\hbar \sum_ip_i\left[\left(\pdiff{t}\ket{\psi_i(t)}\right)\bra{\psi_i(t)}+\ket{\psi_i(t)}\left(\pdiff{t}\bra{\psi_i(t)}\right)\right]\justif{\quad}{Product rule} \\
                                                 & = \sum_ip_i\left[\left(i\hbar\pdiff{t}\ket{\psi_i(t)}\right)\bra{\psi_i(t)}+\ket{\psi_i(t)}\left(-i\hbar\pdiff{t}\bra{\psi_i(t)}\right)\right]     \\
                                                 & = \sum_ip_i\left[\left(\hat{H}\ket{\psi_i(t)}\right)\bra{\psi_i(t)}+\ket{\psi_i(t)}\left(-\bra{\psi_i(t)}\hat{H}\right)\right]                     \\
                                                 & =\hat{H}\left[\sum_ip_i\ket{\psi_i(t)}\bra{\psi_i(t)}\right]-\left[\sum_ip_i\ket{\psi_i(t)}\bra{\psi_i(t)}\right]\hat{H}                           \\
                                                 & =\hat{H}\hat{\rho}-\hat{\rho}\hat{H}=\left[\hat{H},\hat{\rho}(t)\right]
  \end{align*}
\end{soln}
\end{document}